
%%%%%%%%%%%%%%%%%%%%%%%%%%%%%%%%%%%%%%%%%%%%%%%%%%%%%%%%%%%%%
%% Begin exercise %%
%%%%%%%%%%%%%%%%%%%%%%%%%%%%%%%%%%%%%%%%%%%%%%%%%%%%%%%%%%%%%
\ex{Power MOSFETs}

%%%%%%%%%%%%%%%%%%%%%%%%%%%%%%%%%%%%%%%%%%%%%%%%%%%%%%%%%%%%%
%% Task 1 %%
%%%%%%%%%%%%%%%%%%%%%%%%%%%%%%%%%%%%%%%%%%%%%%%%%%%%%%%%%%%%%

\task{Choosing the MOSFET for two converter roles}

Two MOSFETs are available. The corresponding parameters are listed in \autoref{table:mosfet_task1_devices}.

\begin{table}[ht]
    \centering
    \begin{tabular}{lll}
        \toprule
        Parameter & Device A & Device B \\
        \midrule
        Drain-source blocking voltage $V_\mathrm{DSS}$ & \SI{80}{\volt} & \SI{650}{\volt} \\
        On-state resistance $R_\mathrm{DS(on)}$ & \SI{3}{\milli\ohm} & \SI{120}{\milli\ohm} \\
        Switching interval $(t_\mathrm{r} + t_\mathrm{f})$ & \SI{40}{\nano\second} & \SI{100}{\nano\second} \\
        \bottomrule
    \end{tabular}
    \caption{MOSFET data for Task 1.}
    \label{table:mosfet_task1_devices}
\end{table}

Two applications are being considered. \\
\textbf{Application I:} \\
\SI{48}{\volt} synchronous buck leg, load current $I = \SI{40}{\ampere}$, switching frequency $f_\mathrm{s} = \SI{200}{\kilo\hertz}$, \newline
duty ratio $D = 0.5$ \\

\textbf{Application II:} \\
\SI{400}{\volt} boost / PFC switch position, current $I = \SI{8}{\ampere}$, switching frequency $f_\mathrm{s} = \SI{50}{\kilo\hertz}$, \newline
duty ratio $D = 0.5$ \\

Use the slide expressions
\begin{align}
    P_\mathrm{ON} &= I_\mathrm{D(rms)}^2 R_\mathrm{DS(on)} \\
    P_\mathrm{switching} &= \frac{1}{6} V_\mathrm{DS} I_\mathrm{D} (t_\mathrm{r} + t_\mathrm{f}) f_\mathrm{s}.
\end{align}

\subtask{Select the appropriate device for each application.}

\begin{solutionblock}
    The choice is based on the required parameter and the semiconductor characteristic. For application I, both devices are suitable due to their 
    drain-source blocking voltage, on-state resistance and switching frequency. If you compare the single values, device A has lower power losses
    due to lower on-state resistance and lower switching interval. Therefore device A is selected for application I.
    For application II the exclusion criterion for device A is the drain-source blocking voltage of $\SI{80}{\volt}$, but \SI{400}{\volt} is necessary.
    This means for application II you can only choose device B, even if the other parameters of device A are more suitable.
\end{solutionblock}

\subtask{Compute the conduction and switching loss for each candidate in each application.}

\begin{solutionblock}
    For $D = 0.5$, the root-mean-square current is
    \begin{equation}
        I_\mathrm{rms} = I \sqrt{D}.
    \end{equation}

    For Application I:
    \begin{equation}
        I_\mathrm{rms} = \SI{40}{\ampere} \cdot \sqrt{0.5} = \SI{28.3}{\ampere}.
    \end{equation}

    Device A:
    \begin{align}
        P_\mathrm{ON,A} &= \left( \SI{28.3}{\ampere} \right)^2 \cdot \SI{3}{\milli\ohm} \approx \SI{2.4}{\watt} \\
        P_\mathrm{sw,A} &= \frac{1}{6} \cdot \SI{48}{\volt} \cdot \SI{40}{\ampere} \cdot \SI{40}{\nano\second} \cdot \SI{200}{\kilo\hertz}
        \approx \SI{2.56}{\watt} \\
        P_\mathrm{tot,A} &\approx \SI{4.96}{\watt}.
    \end{align}

    Device B:
    \begin{align}
        P_\mathrm{ON,B} &= \left( \SI{28.3}{\ampere} \right)^2 \cdot \SI{120}{\milli\ohm} \approx \SI{96}{\watt} \\
        P_\mathrm{sw,B} &= \frac{1}{6} \cdot \SI{48}{\volt} \cdot \SI{40}{\ampere} \cdot \SI{100}{\nano\second} \cdot \SI{200}{\kilo\hertz}
        \approx \SI{6.4}{\watt} \\
        P_\mathrm{tot,B} &\approx \SI{102.4}{\watt}.
    \end{align}

    Hence, Device A is clearly correct for Application I.

    For Application II:
    \begin{equation}
        I_\mathrm{rms} = \SI{8}{\ampere} \cdot \sqrt{0.5} = \SI{5.66}{\ampere}.
    \end{equation}

    Device B:
    \begin{align}
        P_\mathrm{ON,B} &= \left( \SI{5.66}{\ampere} \right)^2 \cdot \SI{120}{\milli\ohm} \approx \SI{3.84}{\watt} \\
        P_\mathrm{sw,B} &= \frac{1}{6} \cdot \SI{400}{\volt} \cdot \SI{8}{\ampere} \cdot \SI{100}{\nano\second} \cdot \SI{50}{\kilo\hertz}
        \approx \SI{2.67}{\watt} \\
        P_\mathrm{tot,B} &\approx \SI{6.51}{\watt}.
    \end{align}

\end{solutionblock}

\subtask{Explain why the high-voltage MOSFET is penalized in conduction.}

\begin{solutionblock}
    A high-voltage MOSFET needs a thicker and more lightly doped drift region in order to withstand the blocking voltage.
    This raises $R_\mathrm{DS(on)}$ and therefore increases the conduction loss. 
    The device that survives voltage stress is therefore not automatically the device with the lowest on-state loss.
\end{solutionblock}

\subtask{State one design lesson you should extract from this comparison.}

\begin{solutionblock}
    MOSFET selection is not only a matter of blocking voltage. 
    A device with a much higher voltage rating can be a very poor choice at low voltage because its on-state resistance becomes unnecessarily large.
\end{solutionblock}


%%%%%%%%%%%%%%%%%%%%%%%%%%%%%%%%%%%%%%%%%%%%%%%%%%%%%%%%%%%%%
%% Task 2 %%
%%%%%%%%%%%%%%%%%%%%%%%%%%%%%%%%%%%%%%%%%%%%%%%%%%%%%%%%%%%%%

\task{Gate charge, Miller plateau, and switching-speed trade-off}

A MOSFET is driven from a gate driver through a gate resistor $R_\mathrm{g}$ according \autoref{table:MOSFET_Gate_Driver_parameter}.

\begin{table}[ht]
    \centering
    \begin{tabular}{ll}
        \toprule
        Parameter & Value\\
        \midrule
        Driver voltage $V_\mathrm{drv}:$ & $\SI{12}{\volt}$\\
        Threshold voltage $V_\mathrm{GS(th)}:$ & $\SI{4}{\volt}$\\
        Miller plateau voltage $V_\mathrm{GP}:$ & $\SI{6}{\volt}$\\
        Gate resistance $R_\mathrm{g}:$ & $\SI{10}{\ohm}$\\
        Charge from $\SI{0}{\volt}$ to threshold: $Q_\mathrm{GS1}:$ & $\SI{8}{\nano\coulomb}$\\
        Gate-drain charge  $Q_\mathrm{GD}$ during plateau: & $\SI{24}{\nano\coulomb}$\\
        Effective Miller capacitance $C_\mathrm{GD}$ & $\SI{300}{\pico\farad}$\\
        \bottomrule
    \end{tabular}
    \caption{MOSFET and Gate Driver Characteristics.}
    \label{table:MOSFET_Gate_Driver_parameter}
\end{table}

Assume piecewise-constant gate current:
\begin{align*}
    I_\mathrm{g1} &\approx \frac{V_\mathrm{drv}}{R_\mathrm{g}} = \frac{\SI{12}{\volt}}{\SI{10}{\ohm}} = \SI{1.2}{\ampere} \\
    I_\mathrm{g2} &\approx \frac{V_\mathrm{drv} - V_\mathrm{GP}}{R_\mathrm{g}}
    = \frac{\SI{12}{\volt} - \SI{6}{\volt}}{\SI{10}{\ohm}} = \SI{0.6}{\ampere}.
\end{align*}

\subtask{Estimate the delay to reach threshold.}

\begin{solutionblock}
    The delay up to threshold is
    \begin{equation}
        t_\mathrm{d} \approx \frac{Q_\mathrm{GS1}}{I_\mathrm{g1}}
        = \frac{\SI{8}{\nano\coulomb}}{\SI{1.2}{\ampere}}
        \approx \SI{6.7}{\nano\second}.
    \end{equation}
\end{solutionblock}

\subtask{Estimate the Miller plateau duration.}

\begin{solutionblock}
    The Miller interval is estimated by
    \begin{equation}
        t_\mathrm{M} \approx \frac{Q_\mathrm{GD}}{I_\mathrm{g2}}
        = \frac{\SI{24}{\nano\coulomb}}{\SI{0.6}{\ampere}}
        = \SI{40}{\nano\second}.
    \end{equation}
\end{solutionblock}

\subtask{Estimate the voltage slew rate during the Miller plateau using
\begin{equation*}
    \frac{\mathrm{d}V_\mathrm{DS}}{\mathrm{d}t} \approx \frac{I_\mathrm{g}}{C_\mathrm{GD}}.
\end{equation*}
}

\begin{solutionblock}
    During the plateau,
    \begin{equation}
        \frac{\mathrm{d}V_\mathrm{DS}}{\mathrm{d}t}
        \approx \frac{\SI{0.6}{\ampere}}{\SI{300}{\pico\farad}}
        = 2 \cdot 10^9 \,\si{\volt\per\second}
        = \SI{2}{\volt\per\nano\second}.
    \end{equation}
\end{solutionblock}

\subtask{Repeat qualitatively for $R_\mathrm{g} = \SI{20}{\ohm}$ and explain why ``faster'' is not always ``better''.}

\begin{solutionblock}
    If $R_\mathrm{g}$ doubles to \SI{20}{\ohm}, the gate current roughly halves. Therefore,
    the threshold delay roughly doubles, the Miller plateau roughly doubles, and the voltage slew rate roughly halves.
    A larger gate resistance slows the switching transient, usually reduces ringing and EMI, but increases overlap loss.
    Hence, fast switching is not the only requirement because the designer must also consider ringing, overshoot, and electromagnetic interference.
\end{solutionblock}


%%%%%%%%%%%%%%%%%%%%%%%%%%%%%%%%%%%%%%%%%%%%%%%%%%%%%%%%%%%%%
%% Task 3 %%
%%%%%%%%%%%%%%%%%%%%%%%%%%%%%%%%%%%%%%%%%%%%%%%%%%%%%%%%%%%%%

\task{Should the body diode be trusted as the freewheel path?}

A half-bridge uses a MOSFET according \autoref{table:MOSFET_half_bridge_parameter}. The body diode of the MOSFET is conducting 
before the complementary switch turns on.

\begin{table}[ht]
    \centering
    \begin{tabular}{ll}
        \toprule
        Parameter & Value\\
        \midrule
        Bus voltage $V_\mathrm{DC}:$ & $\SI{100}{\volt}$\\
        Reverse-recovery peak current $I_\mathrm{rr}:$ & $\SI{15}{\ampere}$\\
        Reverse-recovery time $t_\mathrm{rr}:$ & $\SI{60}{\nano\second}$   \\
        Switching frequency $f_\mathrm{s}:$ & $\SI{200}{\kilo\hertz}$      \\   
        \bottomrule
    \end{tabular}
    \caption{MOSFET and half-bridge Characteristics.}
    \label{table:MOSFET_half_bridge_parameter}
\end{table}

Approximate:
\begin{align*}
    Q_\mathrm{rr} &\approx \frac{1}{2} I_\mathrm{rr} t_\mathrm{rr} \\
    E_\mathrm{rr} &\approx \frac{1}{2} Q_\mathrm{rr} V_\mathrm{DC}.
\end{align*}

\subtask{Estimate the reverse-recovery charge $Q_\mathrm{rr}$.}

\begin{solutionblock}
    The reverse-recovery charge $Q_\mathrm{rr}$ is calculated by
    \begin{equation}
        Q_\mathrm{rr}
        = \frac{1}{2} \cdot \SI{15}{\ampere} \cdot \SI{60}{\nano\second}
        = \SI{450}{\nano\coulomb}
        = \SI{0.45}{\micro\coulomb}.
    \end{equation}
\end{solutionblock}

\subtask{Estimate the reverse-recovery energy $E_\mathrm{rr}$.}

\begin{solutionblock}
    The reverse-recovery energy $E_\mathrm{rr}$ is obtained by
    \begin{equation}
        E_\mathrm{rr}
        = \frac{1}{2} \cdot \SI{0.45}{\micro\coulomb} \cdot \SI{100}{\volt}
        = \SI{22.5}{\micro\joule}.
    \end{equation}
\end{solutionblock}

\subtask{Estimate the average recovery power.}

\begin{solutionblock}
    The average recovery power yields
    \begin{equation}
        P_\mathrm{rr} = E_\mathrm{rr} f_\mathrm{s}
        = \SI{22.5}{\micro\joule} \cdot \SI{200}{\kilo\hertz}
        = \SI{4.5}{\watt}.
    \end{equation}
\end{solutionblock}

\subtask{Explain why a fast voltage rise at the opposite switch can become dangerous, and state when an external freewheeling diode becomes attractive.}

\begin{solutionblock}
    The body diode alone may already produce a non-trivial reverse-recovery loss. 
    More importantly, a very fast reapplied drain-source voltage plus reverse-recovery current can interact with the parasitic bipolar transistor and damage the MOSFET.

    An external freewheeling diode becomes attractive when the commutation is hard-switched, the switching frequency is high, and the body diode recovery becomes a relevant source of loss or stress.
\end{solutionblock}


%%%%%%%%%%%%%%%%%%%%%%%%%%%%%%%%%%%%%%%%%%%%%%%%%%%%%%%%%%%%%
%% Task 4 %%
%%%%%%%%%%%%%%%%%%%%%%%%%%%%%%%%%%%%%%%%%%%%%%%%%%%%%%%%%%%%%

\task{Parallel MOSFETs: why current sharing is naturally stabilizing}

Two identical MOSFETs are connected in parallel. At \SI{25}{\celsius},
\begin{equation}
    R_\mathrm{DS(on),25} = \SI{4}{\milli\ohm}.
\end{equation}

Assume the temperature coefficient
\begin{equation}
    R_\mathrm{DS(on)}(T) = R_{25} \left[ 1 + 0.004(T - 25) \right].
\end{equation}

At one operating point:
\begin{itemize}
    \setlength{\itemsep}{-2pt} % Distance between the points (Part1)
    \setlength{\parskip}{-2pt} % Distance between the points (Part2)    
    \item MOSFET 1 is at \SI{100}{\celsius}
    \item MOSFET 2 is at \SI{50}{\celsius}
    \item the pair carries a total current of \SI{120}{\ampere}
\end{itemize}

\subtask{Calculate the two on-state resistances.}

\begin{solutionblock}
    For MOSFET 1 the on-state resistances is calculated by
    \begin{equation}
        R_1 = \SI{4}{\milli\ohm} \cdot \left[ 1 + 0.004(100 - 25) \right]
        = \SI{5.2}{\milli\ohm}.
    \end{equation}

    For MOSFET 2 the on-state resistances is calculated by
    \begin{equation}
        R_2 = \SI{4}{\milli\ohm} \cdot \left[ 1 + 0.004(50 - 25) \right]
        = \SI{4.4}{\milli\ohm}.
    \end{equation}
\end{solutionblock}

\subtask{Assuming equal drain-source voltage, find the current split.}

\begin{solutionblock}
    The current splits inversely with resistance:
    \begin{align}
        I_1 &= \SI{120}{\ampere} \cdot \frac{R_2}{R_1 + R_2}
        = \SI{120}{\ampere} \cdot \frac{4.4}{5.2 + 4.4}
        = \SI{55}{\ampere} \\
        I_2 &= \SI{120}{\ampere} \cdot \frac{R_1}{R_1 + R_2}
        = \SI{120}{\ampere} \cdot \frac{5.2}{5.2 + 4.4}
        = \SI{65}{\ampere}.
    \end{align}
\end{solutionblock}

\subtask{Explain why this tends to stabilize sharing and contrast it with BJT paralleling.}

\begin{solutionblock}
    The hotter MOSFET has the larger on-state resistance and therefore takes less current. 
    This negative feedback makes current sharing naturally stabilizing.

    In contrast, BJTs do not show this same helpful positive temperature coefficient in the on-state and are much more prone to current hogging.
\end{solutionblock}

